%% bare_jrnl_compsoc.tex
%% V1.4b
%% 2015/08/26
%% by Michael Shell
%% See:
%% http://www.michaelshell.org/
%% for current contact information.
%%
%% This is a skeleton file demonstrating the use of IEEEtran.cls
%% (requires IEEEtran.cls version 1.8b or later) with an IEEE
%% Computer Society journal paper.
%%
%% Support sites:
%% http://www.michaelshell.org/tex/ieeetran/
%% http://www.ctan.org/pkg/ieeetran
%% and
%% http://www.ieee.org/

%%*************************************************************************
%% Legal Notice:
%% This code is offered as-is without any warranty either expressed or
%% implied; without even the implied warranty of MERCHANTABILITY or
%% FITNESS FOR A PARTICULAR PURPOSE! 
%% User assumes all risk.
%% In no event shall the IEEE or any contributor to this code be liable for
%% any damages or losses, including, but not limited to, incidental,
%% consequential, or any other damages, resulting from the use or misuse
%% of any information contained here.
%%
%% All comments are the opinions of their respective authors and are not
%% necessarily endorsed by the IEEE.
%%
%% This work is distributed under the LaTeX Project Public License (LPPL)
%% ( http://www.latex-project.org/ ) version 1.3, and may be freely used,
%% distributed and modified. A copy of the LPPL, version 1.3, is included
%% in the base LaTeX documentation of all distributions of LaTeX released
%% 2003/12/01 or later.
%% Retain all contribution notices and credits.
%% ** Modified files should be clearly indicated as such, including  **
%% ** renaming them and changing author support contact information. **
%%*************************************************************************


% *** Authors should verify (and, if needed, correct) their LaTeX system  ***
% *** with the testflow diagnostic prior to trusting their LaTeX platform ***
% *** with production work. The IEEE's font choices and paper sizes can   ***
% *** trigger bugs that do not appear when using other class files.       ***                          ***
% The testflow support page is at:
% http://www.michaelshell.org/tex/testflow/


\documentclass[10pt,journal,compsoc]{IEEEtran}

% *** CITATION PACKAGES ***
%
\ifCLASSOPTIONcompsoc
  % IEEE Computer Society needs nocompress option
  % requires cite.sty v4.0 or later (November 2003)
  \usepackage[nocompress]{cite}
\else
  % normal IEEE
  \usepackage{cite}
\fi

\hyphenation{op-tical net-works semi-conduc-tor}

\newcommand*{\sectiondir}{sections/}

\begin{document}

% EDIT TO CHANGE THE TITLE
\title{Demo for Data and Network Security report}
%
%
% author names and IEEE memberships
% note positions of commas and nonbreaking spaces ( ~ ) LaTeX will not break
% a structure at a ~ so this keeps an author's name from being broken across
% two lines.
% use \thanks{} to gain access to the first footnote area
% a separate \thanks must be used for each paragraph as LaTeX2e's \thanks
% was not built to handle multiple paragraphs
%
%
%\IEEEcompsocitemizethanks is a special \thanks that produces the bulleted
% lists the Computer Society journals use for "first footnote" author
% affiliations. Use \IEEEcompsocthanksitem which works much like \item
% for each affiliation group. When not in compsoc mode,
% \IEEEcompsocitemizethanks becomes like \thanks and
% \IEEEcompsocthanksitem becomes a line break with idention. This
% facilitates dual compilation, although admittedly the differences in the
% desired content of \author between the different types of papers makes a
% one-size-fits-all approach a daunting prospect. For instance, compsoc 
% journal papers have the author affiliations above the "Manuscript
% received ..."  text while in non-compsoc journals this is reversed. Sigh.

% EDIT HERE TO ADD THE AUTHORS
\author{Name~Surname,~\IEEEmembership{1000000,~Sapienza University of Rome,}\\
        Name~Surname,~\IEEEmembership{1000000,~Sapienza University of Rome,}%


% EDIT HERE TO ADD THE DELIVERY DATE
\thanks{Manuscript delivered on Month Day Year}}


% EDIT HERE TO ADD THE ABSTRACT
\IEEEtitleabstractindextext{%
\begin{abstract}
The abstract goes here.
\end{abstract}

% EDIT TO ADD THE KEYWORDS OF YOUR REPORT (e.g., machine learning, adversarial attacks)
\begin{IEEEkeywords}
Machine Learning, poisoning attacks, watermarking...
\end{IEEEkeywords}}


% make the title area
\maketitle


\IEEEdisplaynontitleabstractindextext
\IEEEpeerreviewmaketitle

\section{Introduction}\label{sec:introduction}

In this section you need to introduce the topic, explain the problem you want to tackle and why it is something interesting to treat, 
the contribution you want to propose and why your proposal is better than current state-of-the-art. Is the project academic or industrially oriented?
\section{Background}\label{sec:background}

In this section you need to give the reader background information about the topic, 
e.g., if you are talking about deep neural networks, explain a bit how they work. What is in the marketplace or academic literature?
\section{Overview of your proposed approach}\label{sec:our_approach}

In this section you need to provide scientific and technical detail about your proposal, explaining what you want to do, describing the idea and how it differs to current state-of-the-art. How can your project be done?
Select approaches to form the basis of the solution, ensure you can justify it
\section{Evaluation}\label{sec:evaluation}

In this section you should propose an experiment to evaluate the goodness of your proposal. Of course for this report you are NOT REQUIRED TO DO ANY EXPERIMENTS, it's just a PROPOSAL of experiments that may be useful to do.
\section{Related work}\label{sec:related_work}

In this section, you need to analyze your direct competitors (e.g.~\cite{adi_watermarking} if you talk about watermarking), explaining a bit what they do and why you think your approach is better. What has been tried before? 
\section{Conclusions}\label{sec:conclusions}

This section concludes the report providing a brief summary of what can be accomplished within the project and possible further future extensions.

PS: It is mandatory to include the section with bibliographic references in this research report


\bibliographystyle{IEEEtranS}
\bibliography{bibliography}


% that's all folks
\end{document}


